\documentclass[11pt,a4paper]{article}

% ========================
% ESSENTIAL PACKAGES
% ========================

% Language and encoding
\usepackage[utf8]{inputenc}
\usepackage[T1]{fontenc}
\usepackage[english]{babel}
\usepackage{natbib}

% Math packages
\usepackage{amsmath}        % Enhanced math environments
\usepackage{amssymb}        % Additional math symbols
\usepackage{amsthm}         % Theorem environments
\usepackage{amsfonts}       % Additional math fonts
\usepackage{mathtools}      % Extension to amsmath
\usepackage{bm}             % Bold math symbols
\usepackage{dsfont}         % Double-struck fonts (for sets like ℝ, ℕ)
\usepackage{mathrsfs}       % Script fonts for math

% Graphics and figures
\usepackage{graphicx}       % Include graphics
\usepackage{float}          % Better float placement
\usepackage{subcaption}     % Subfigures
\usepackage{tikz}           % Drawing diagrams
\usepackage{pgfplots}       % Plotting
\pgfplotsset{compat=1.18}

% Page layout and formatting
\usepackage[margin=1in]{geometry}
\usepackage{fancyhdr}       % Custom headers and footers
\usepackage{titlesec}       % Customize section titles
\usepackage{enumitem}       % Customize lists
\usepackage{parskip}        % Paragraph spacing

% Colors and highlighting
\usepackage{xcolor}
\usepackage{mdframed}       % Framed environments
\usepackage{tcolorbox}      % Colored boxes

% Tables
\usepackage{array}          % Extended array environments
\usepackage{booktabs}       % Professional tables
\usepackage{multirow}       % Multi-row cells

% References and links
\usepackage{hyperref}       % Hyperlinks
\usepackage{cleveref}       % Smart cross-references

% Additional useful packages
\usepackage{cancel}         % Cancel terms in equations
\usepackage{siunitx}        % SI units
\usepackage{listings}       % Code listings (if needed)

% ========================
% CUSTOM COMMANDS
% ========================

% Common math sets
\newcommand{\R}{\mathbb{R}}
\newcommand{\N}{\mathbb{N}}
\newcommand{\Z}{\mathbb{Z}}
\newcommand{\Q}{\mathbb{Q}}
\newcommand{\C}{\mathbb{C}}

% Common operators
\DeclareMathOperator{\lcm}{lcm}
\DeclareMathOperator{\Tr}{Tr}
\DeclareMathOperator{\rank}{rank}
%\DeclareMathOperator{\span}{span}
%\DeclareMathOperator{\null}{null}

% Shortcuts for common expressions
\newcommand{\abs}[1]{\left|#1\right|}
\newcommand{\norm}[1]{\left\|#1\right\|}
\newcommand{\inner}[2]{\left\langle #1, #2 \right\rangle}
\newcommand{\ceil}[1]{\left\lceil #1 \right\rceil}
\newcommand{\floor}[1]{\left\lfloor #1 \right\rfloor}

% Derivatives and integrals
\newcommand{\dd}[1]{\,\mathrm{d}#1}
\newcommand{\dv}[2]{\frac{\mathrm{d}#1}{\mathrm{d}#2}}
\newcommand{\pdv}[2]{\frac{\partial #1}{\partial #2}}

% ========================
% THEOREM ENVIRONMENTS
% ========================

\theoremstyle{definition}
\newtheorem{definition}{Definition}[section]
\newtheorem{example}{Example}[section]
\newtheorem{exercise}{Exercise}[section]

\theoremstyle{plain}
\newtheorem{theorem}{Theorem}[section]
\newtheorem{lemma}{Lemma}[section]
\newtheorem{corollary}{Corollary}[section]
\newtheorem{proposition}{Proposition}[section]

\theoremstyle{remark}
\newtheorem{remark}{Remark}[section]
\newtheorem{note}{Note}[section]

% ========================
% COLORED BOXES
% ========================

% Important theorem box
\newtcolorbox{importantbox}{
	colback=blue!5!white,
	colframe=blue!75!black,
	title=Important
}

% Warning/Note box
\newtcolorbox{notebox}{
	colback=yellow!10!white,
	colframe=orange!75!black,
	title=Note
}

% Example box
\newtcolorbox{examplebox}{
	colback=green!5!white,
	colframe=green!75!black,
	title=Example
}

% ========================
% DOCUMENT SETUP
% ========================

% Page style
\pagestyle{fancy}
\fancyhf{}
\rhead{\thepage}
\lhead{\leftmark}

% Title formatting
\title{Functional Bandit Notes}
\author{Your Name}
\date{\today}

% Hyperlink setup
\hypersetup{
	colorlinks=true,
	linkcolor=blue,
	filecolor=magenta,      
	urlcolor=cyan,
	citecolor=red
}

% ========================
% DOCUMENT CONTENT
% ========================

\begin{document}
	
\section{Model}
Consider the RKHS framework proposed by \cite{cai2012minimax}.
\begin{itemize}
 \item Consider two-armed bandit with i.i.d. square integrable functional contexts \(x_1,x_2,x_3,\ldots,x_T\).
 \item There exists $\beta_1,\beta_2\in \mathcal{H}_K$ which is an RKHS with kernel $K,$ and the rewards for two arms are
 \[f^{(1)}(x) = \langle x,\beta_1\rangle+\varepsilon \quad\text{and}\quad f^{(1)}(x) = \langle x,\beta_2\rangle+\varepsilon. \]
 \item In the offline setting, we assume that the operator $L_{K^{1/2} C K^{1/2}}$ has a decaying spectrum, where $C$ is the covariance function of $x.$
\end{itemize}
\section{Algorithn}
\subsection{Successive elimination}
	We need to find $\Delta(x)$ such that with high probability, $|\langle x_t, \hat{\beta}-\beta\rangle|\leq \Delta(x)$ for every $t\in [T].$ We also need to bound the expectation of $\Delta(x)$ on the population of $x.$

By \cite{cai2012minimax}, we have
\begin{align*}
	\langle x_t, \hat{\beta}-\beta\rangle = \langle x_t, L_{K^{1/2}} (\hat{f}_{\lambda}-f)\rangle
=\langle x_t, L_{K^{1/2}} (\hat{f}_{\lambda}-f_{\lambda})\rangle
+\langle x_t, L_{K^{1/2}} ({f}_{\lambda}-f_0)\rangle,
\end{align*}
where the first term on the RHS is the variance term and the second term on the RHS is the bias term.
The bias term is deterministic given $x.$ In \cite{cai2012minimax}, the bias term is handled by the interplay between  the covariance of $x_t$ and $f_\lambda-f_0.$ Basically, $f_\lambda-f_0$ can be decomposited by the eigenfunctions of $L_{K^{1/2}CK^{1/2}}.$ 
%	\section{References}
%	
%	You can reference equations \eqref{eq:example}, theorems \cref{thm:example}, and figures \cref{fig:parabola} using cleveref.
	
	% Add bibliography if needed
\bibliographystyle{plainnat}
\bibliography{ref}
	
\end{document}